\section*{第 \ref{ch:b3:extension-strike-slip} 章——伸展与走滑构造}

\begin{solution}{exo:b3:extension-strike-slip:1}
岩层倾角为 $50 - 35 = \ang{15}$，倾向远离断层；$\beta =
\sin\ang{50}/\sin\ang{35} = 0.766/0.574 = 1.34$。
\end{solution}

\begin{solution}{exo:b3:extension-strike-slip:2}
右行错列向右：盆地。右行错列向左：脊。左行
错列向左：盆地。左行错列向右：脊。
\end{solution}

\begin{solution}{exo:b3:extension-strike-slip:3}
$105/20 = \qty{5.25}{km/Ma} = \qty{5.25}{mm/\text{年}}$，与
\qty{6}{mm/\text{年}} 接近：在 \qty{20}{Ma} 内速率不必保持恒定。
\end{solution}

\begin{solution}{exo:b3:extension-strike-slip:4}
\qty{200}{\percent} 表示 $\Delta L = 2L_0$：$L_0 = \qty{117.5}{km}$，$L
= \qty{352.5}{km}$。在北部，$\Delta L = 0.5L_0$：$L_0 = \qty{470}{km}$，
和 $L = \qty{705}{km}$。
\end{solution}

\begin{solution}{exo:b3:extension-strike-slip:5}
$563/20 = 28$ 至 $563/15 = \qty{38}{mm/\text{年}}$。该位移量为最小值，而
跨断裂带测量得到的漂移还包括
该断层系中邻近断层的运动，以及
地震间积累的弹性应变。
\end{solution}

\begin{solution}{exo:b3:extension-strike-slip:6}
$t = 90 - \text{倾角}$。在 \ang{60} 处：$(1 + 0.6\times1.732)/(1 - 0.6\times
0.577) = 2.039/0.654 = 3.12$；在 \ang{50} 处：$1.715/0.497 = 3.45$；在
\ang{40} 处：$1.504/0.285 = 5.28$。锁定角为 $\arctan0.6 = \ang{31.0}$。
\end{solution}

\begin{solution}{exo:b3:extension-strike-slip:7}
$\arctan0.4 = \ang{21.8}$，$\arctan0.2 = \ang{11.3}$。倾角为
\ang{15} 的断层仅当 $\mu < \tan\ang{15} = 0.27$ 时才会滑动。
\end{solution}

\begin{solution}{exo:b3:extension-strike-slip:8}
$2.5\times4.0 = \qty{10.0}{km^2}$，$2.5\times(4.0 - 1.5) =
\qty{6.25}{km^2}$。
\end{solution}

\begin{solution}{exo:b3:extension-strike-slip:9}
$132/1.19 = \qty{110.9}{km}$，伸展了 \qty{21.1}{km}，这需要
$21.1/6 = \qty{3.5}{Ma}$。
\end{solution}

\begin{solution}{exo:b3:extension-strike-slip:10}
之前为 $3.0/\sin\ang{60} = \qty{3.464}{km}$，之后为 $3.0/\sin\ang{45} =
\qty{4.243}{km}$：每个块体 \qty{0.779}{km}（$\beta = 1.22$）。岩层
倾角为 $60 - 45 = \ang{15}$。每个块体整体转动：
使向东倾断层的倾角变小的转动，会使块体中的水平线
向东抬升，岩层向西倾，与
断层的倾向相反。
\end{solution}

\begin{solution}{exo:b3:extension-strike-slip:11}
在左行东西向带中，$\sigma_1$ 走向为 \ang{045}。R 剪切面
与带成 $\phi/2 = \ang{20}$，位于 $\sigma_1$ 一侧：走向 \ang{070}，左行。
R$'$ 与带成 $90 - 20 = \ang{70}$：走向 \ang{020}，右行。张裂隙
平行于 $\sigma_1$：走向 \ang{045}。
\end{solution}

\begin{solution}{exo:b3:extension-strike-slip:12}
当 $\mu = 0.6$ 时，断层形成于 \ang{60.5}，锁定于 \ang{31.0}：$\beta =
\sin\ang{60.5}/\sin\ang{31.0} = 0.870/0.515 = 1.69$。当 $\mu = 0.85$ 时，由
\ang{65.2} 至 \ang{40.4}，$\beta = 0.908/0.648 = 1.40$。摩擦力较小的断层
在锁定前转动更多：弱岩石在
必须形成新断层之前伸展得更多。
\end{solution}

\begin{solution}{pb:b3:extension-strike-slip:1}
\textbf{1.} 岩层原为水平，并随块体一同转动；
断层转过了相同的角度，因此 $\theta_0 = 42 + 18 = \ang{60}$，
即 $\phi = \ang{30}$ 时的安德森倾角。

\textbf{2.} $\beta = \sin\ang{60}/\sin\ang{42} = 0.8660/0.6691 = 1.294$。

\textbf{3.} $24.0/1.294 = \qty{18.5}{km}$；该区域已伸展了
\qty{5.5}{km}（\qty{5.46}{km}）。

\textbf{4.} $t = \ang{48}$：$(1 + 0.6\times0.900)/(1 - 0.6\times1.111) =
1.540/0.334 = 4.62$，与新断层在 \ang{60.5} 处的 $3.12$ 相比，
多出约一半。

\textbf{5.} 在 $\arctan0.6 = \ang{31.0}$ 处；$\beta = \sin\ang{60}/\sin\ang{31}
= 1.68$ 为最大值。在此之前，约 \ang{60} 的新断层会切
割块体，并由第二代多米诺构造接替。

\textbf{6.} 上盘的每一列岩块，其下降量等于其下方断层
在水平断距上的下降量；断层在近地表处最陡，
远离地表则变缓，因此靠近断层的岩块下降最多，远离
断层的岩块下降最少：岩层向断层倾斜。

\textbf{7.} $D = A/H = 16.0/3.2 = \qty{5.0}{km}$。

\textbf{8.} $D$ 与 $A$ 成正比：$\qty{5.0}{km} \pm \qty{0.5}{km}$。

\textbf{9.} $3.2\times(5.0 - 2.0) = \qty{9.6}{km^2}$。

\textbf{10.} 水平断距，\qty{3.2}{km}。

\textbf{11.} 充填物向断层方向增厚，呈楔形，较老的
地层比较年轻的地层倾斜更大：每一层沉积时，
转折端与断层旁的空间都在持续扩大。

\textbf{12.} $z = P/(\rho g) = \num{5.5e8}/(2800\times9.8) = \qty{20.0}{km}$。

\textbf{13.} $20.0/(24 - 16) = \qty{2.5}{km/Ma} = \qty{2.5}{mm/\text{年}}$。

\textbf{14.} $20.0/\sin\ang{30} = \qty{40}{km}$ 为滑距，水平伸展为 $40\cos\ang{30} =
\qty{34.6}{km}$。

\textbf{15.} 否：\ang{15} 低于锁定倾角 \ang{31}。可能是断层先在较陡时滑动，
随后发生转动（\omterm{def:b3:extension-strike-slip:core-complex}{滚动枢纽}），或者
断层较弱，$\mu < 0.27$，由黏土或高压流体所致。

\textbf{16.} 由 $45 - 15 = \ang{30}$ 可知：上盘远离时，
失去负荷的下盘通过地壳均衡抬升并弯曲，使断层绕
转折端发生转动。

\textbf{17.} \qty{34.6}{km} 对比半地堑的 \qty{8.7}{km}（第
23 题）：约为其四倍。核心杂岩位于伸展最强处，如
盆岭省中部。

\textbf{18.} 左行断层向左阶步：形成\omterm{def:b3:extension-strike-slip:bends}{释压性弯曲}，盆地张开。

\textbf{19.} $\qty{5}{mm/\text{年}}\times\qty{1.2}{Ma} = \qty{6.0}{km}$。

\textbf{20.} 面积为 $4.0\times6.0 = \qty{24}{km^2}$；长为 $3.0 + 6.0 = \qty{9.0}{km}$，
长宽比为 $9.0/4.0 = 2.25$。

\textbf{21.} $132/10 = 13$ 至 $132/7 = 19$：死海盆地远为
狭长，其位移是宽度的数倍。

\textbf{22.} \omterm{def:b3:extension-strike-slip:bends}{限制性弯曲}：两支之间隆起一脊；剖面上为正\omterm{def:b3:extension-strike-slip:flower}{花状构造}，
断层向上分叉为若干支，各支带有逆冲分量，
类似死海转换断层上宽 100 至 \qty{300}{m} 的
那一处弯曲。

\textbf{23.} $3.2 + 5.46 = \qty{8.66}{km}$：梅罗半地堑
已伸展 \textbf{\qty[text-series-to-math=true]{8.7}{km}}。
\end{solution}
